\section{Introduction}\label{sec:v14-introduction}
A causal simulation has two outputs: the report delivered to its consumer
and the state that must survive until the next report. Classical comparison
records the first and may leave the second unpriced. We keep both. The
quantity of interest is the least error attainable by one task- and
parameter-independent simulator at a prescribed persistent-state budget.
The error compares the entire feedback transcript together with its actual
latent mark. In particular, equality of report marginals does not replace
equality of a marked experiment.

This distinction leads to three optimization problems. A private simulator
retains only its declared state and uses fresh randomness. A hidden-selector
simulator first selects a policy but does not reveal that choice to the
consumer. A visible-selector simulator reveals it and includes it in the
compared joint interface. Hidden convexification has a finite testing dual,
but its realization may require storing a selector index. At a fixed
private width this is not a harmless convention: a delayed three-symbol
channel has zero hidden-selector deficiency and strictly positive private
deficiency at width two. If the selector is visible, its optimized zero set
is again the private zero set.

\subsection{The main statement}
For finite marked causal experiments $E,F$, prescribed widths
$K=(K_0,\ldots,K_T)$ and a fixed finite seed signature, let
$\delta_K(E,F)$ be the infimum of the marked feedback error over all
admissible simulators. The definitions and exact signature-dependent
variational formulas are given in Section~\ref{sec:v13-deficiency}.
The central result can be stated as follows.

\begin{quote}
At each fixed private or finite-selector budget, intrinsic causal deficiency
has a complete family of finite local testing certificates and original-budget
upper witnesses. These bounds converge with an explicit causal modulus.
They compose across stateless marked presentation changes without increasing
the budget, and active microscopic approximation supplies such changes with
finite graph-residual bounds.
\end{quote}

More precisely, Theorem~\ref{thm:v14-local} constructs rational lower and
upper bounds $L_M\le\delta_K\le U_M$ from clipped-simplex policy cells.
In the private case, with at most $d_t$ joint output--state labels in a
row at time $t$, their gap is at most
\[
 1-\prod_{t=0}^T\bigl(1-\min\{1,d_t/(2M)\}\bigr),
\]
with a zero contribution from one-point rows. The modulus depends on depth
and row alphabets, not on the number of histories. Enumeration of the
cells can nevertheless be large; no polynomial-time synthesis conclusion
is drawn. The lower bound assigns a marked feedback test to each cell.
The upper bound is a single admissible row table at the original budget.
This is a nonlinear testing converse, not an unpriced application of
convex separation. For dyadic $M$ the bounds are monotone and converge.
Strict feasible and infeasible thresholds have finite witnesses; exact
boundary decisions use classical real algebraic elimination.

The optimized spectra have the composition law
\[
 \delta_{KL}(E,G)\le\delta_K(E,F)+\delta_L(F,G),
\]
with selector costs multiplied when present. Theorem~\ref{thm:v14-stability}
then transfers an entire intrinsic interval across two-sided stateless
presentation changes. Its lower statement excludes all original-budget
simulators, even when the original carriers are nonfinite. Its upper
statement constructs a simulator by composing the rational row table with
the bridge channels. Thus the scope of an assertion is visible in both
its error and its persistent resources.

The other parts of the article give operational and analytic realizations
of this comparison. Theorem~\ref{thm:v13-endogenous} identifies the exact
common-encoder face for endogenously controlled finite tasks by reachable
compatible covers, and relates the lossy problem to Bellman regrets.
Theorem~\ref{thm:v13-average} treats nonsummable average criteria for
stationary finite-register controllers with actual simultaneous physical
and register resets; deterministic or singular physical transitions are
permitted between resets. Theorem~\ref{thm:v13-active} supplies active
hard-sphere interventions and convergent graph-core approximations.
Theorem~\ref{thm:v14-residual} strengthens their finite-horizon comparison
by an a posteriori bound involving only graph and intervention Gram data.
Its root-sum-square word estimate retains the actual mark and precedes
optimization over every finite consumer width. A finite age cutoff adds
an explicit geometric remainder for average risks.

\subsection{Comparison with classical results}
Blackwell--Le Cam comparison, convex minimax, and finite-dimensional convex
representation are classical ingredients \cite{Blackwell1953,LeCam1964}.
Filtered comparison is not new as a subject. Weisshaupt's filtered
stochastic-operator compactness lemma and randomization theorem
\cite[Lemma 1 and Theorem 5]{Weisshaupt2002} give an explicit nearby
convex comparison result with a finitely additive target space. They do
not impose the persistent-state budget considered here. We distinguish
that retrieved statement from Norberg's closely related filtered
comparison paper \cite{Norberg2002}; the full original-source comparison
with the latter is not claimed complete.

The common-information dynamic program of
Nayyar--Mahajan--Teneketzis \cite[Theorem 3 and Corollary 5]{NayyarMahajanTeneketzis2013}
provides a classical coordinator formulation, including a finite-memory
specialization. The compatible-cover device also has classical
incomplete-machine antecedents \cite{PaullUnger1959}. Our exact task theorem
adds the common encoder, policy-dependent reachable support and Bellman
regret conditions; the cover terminology by itself is not a novelty claim.
Real quantifier elimination is used only for an exact decision corollary
\cite{Basu2014}. Similarly, graph-core semigroup approximation and the
well-defined microscopic hard-sphere flow are credited to their classical
sources \cite{Trotter1958,GallagherSaintRaymondTexier}. The residual proof below uses a
finite-dimensional variation-of-constants formula and does not propose a
new abstract semigroup approximation theorem.

The mathematical distinction retained after those ingredients are removed
is the simultaneous original-budget testing and realization statement:
local feedback tests certify the private feasible set that a convex
filtered comparison can enlarge, and the same optimized quantity accepts
active microscopic error bounds without a hidden consumer-memory charge.
The finite delayed-channel separation makes that distinction explicit.
It is not a claim that all ingredients, or their entire conjunction, have
an exhaustively certified priority. The source-level literature audit
records which original theorems were inspected and which remain to be
compared.

\subsection{Organization and scope}
Section~\ref{sec:v13-deficiency} develops the signature-indexed deficiency,
its local certificates and its composition-stability theorem. The next
sections treat endogenous tasks, regenerative average criteria, and active
microscopic realization. All five sections include the proofs they use.
The complete development retains the preceding mathematical bodies,
including dominated finite-prefix and summable infinite-time duality,
marked continuation quotients, exact synthesis, operator memory, Gaussian
resource laws, and the historical model interfaces.

The physical results fix particle number, action alphabet, observation
noise, and the initial-mark channel. They do not imply the historical
nonlinear kinetic action-sublevel corrector or a Boltzmann--Grad limit.
Likewise an intrinsic lower certificate is a mathematical inequality over
a specified finite model, not an assertion that floating-point Gram
entries certify a physical system. The eleven-component program is a
mathematical dependency graph with adapters and independent lines;
repository chronology and its separately frozen audit are not hypotheses
of any theorem in this article.
